% Chapter 7 -- The Desktop and the Finder (Part II, Reference).
% Sources (Law 1): os/flair/finder_desktop.c / .h (the two desktop icons:
% the volume, labelled with the disk's volume label, top right; the Trash
% below it, fixed; 32 x 32 icons with Geneva 9 labels; the Trash icon full
% or empty); os/milton/kmain.c (DESKTOP.DB: desktop icon places and the
% views of the Finder windows, saved on a move, a close and the power
% commands; the icons in a window laid out from the grid on every open;
% the click that makes the Finder foreground); os/flair/finder_windows.h /
% .c (window titles: "Drive A Files" for the top level, the 8.3 name of a
% folder, "Trash"; the cascade from (20,60) by 20 px; FINDER_WIN_MAX 4;
% 64 icons a window; the 68 x 52 grid in directory order; the hidden
% TRASH, DESKTOP.DB and volume label; .EXE is the application icon);
% os/milton/desktop_db.c (the DESKTOP.DB file).  Model: System 6 guide Ch.
% 3 "Working With the Finder Desktop" and Mac User's Guide 1991 Ch. 10
% "Organizing Your Files" (apple-macintosh.md Sec 8.1, 8.2).
% GATES: test-finder-desktop test-finder-windows test-desktop-db test-flair-desktop-icons test-flair-disk-windows test-flair-finder-service test-flair-file-ops test-pubs-flair-screens
\Gates{test-finder-desktop, test-finder-windows, test-desktop-db,
  test-flair-desktop-icons, test-flair-disk-windows,
  test-flair-finder-service, test-flair-file-ops, test-pubs-flair-screens}
\chapter{The Desktop and the Finder}
\label{ch:finder}

The \term{Finder} is the part of FLAIR that shows you your disk, folders and
files and lets you organize them.  It is always running.  It owns the
desktop, the icons on it, and the windows that show the contents of the
disk and its folders, and it is the active application whenever you click
the desktop or one of its windows.

\section{The desktop}
\index{desktop}
The desktop holds two icons.

\SideHead{The disk icon}At the top right, the icon of the data disk,
labelled with the disk's name, \texttt{INITECH}.  Double-click it to open
the window of the disk's top level.  You can drag it anywhere on the
desktop; and when you drop a file or folder on it, the file or folder moves
to the top level of the disk.

\SideHead{The Trash}At the bottom right, the Trash.  Drop a file or folder
on it to throw it away (Chapter~\ref{ch:trash}); double-click it to open the Trash
window.  The Trash icon bulges while the Trash holds something.  It cannot
be moved.

\FlairScreen[label=fig:icondrag]{icondrag}{The disk icon, moved}

FLAIR records the place of each desktop icon, and the place, size and
state of each Finder window, in a file on the data disk named
\file{DESKTOP.DB}.  It writes the file when you move an icon on the
desktop, when you move or close a Finder window, and when you choose
\menuitem{Restart}, \menuitem{Shut Down} or \menuitem{Exit to InitechDOS}.
If \file{DESKTOP.DB} or the Trash folder is missing when FLAIR starts,
FLAIR makes a new one.

\section{Finder windows}
\index{Finder windows}\index{windows!Finder}
Each Finder window shows the contents of one folder of the disk.

\begin{center}\small
\begin{tabular}{@{}ll@{}}
\textbf{Window title} & \textbf{Shows}\\\hline
\texttt{Drive A Files} & the top level of the data disk\\
the folder's name, such as \texttt{APPS} & that folder\\
\texttt{Trash} & the Trash\\
\end{tabular}
\end{center}

The first window opens with its top left corner just below the menu bars at
the left of the screen; each window you open after it is placed a little
lower and further right, so that their title bars can all be seen.  A
folder has only one window: if its window is already open, opening the
folder again brings that window to the front.

Up to four Finder windows can be open at a time, and a window shows up to
64 icons.  The icons are laid out in rows from the top left, in the order
in which the files are stored in the folder, each time the window opens;
you can move them within the window while it is open, and tidy them with
\menuitem{Clean Up} or \menuitem{Arrange (by Name)}.

\section{Icons}
\index{icons}
\FlairScreen[label=fig:iconkinds]{iconkinds}{Applications and documents in
the APPS folder}

\SideHead{Applications}Files whose names end in \file{.EXE} are shown as
a diamond.  Double-click one to open the application (Chapter~\ref{ch:applications}).

\SideHead{Documents and programs}Every other file is shown as a page with
a folded corner.  These files are not opened from the desktop: the
programs among them, whose names end in \file{.COM}, run at the InitechDOS
command prompt, and documents are opened from within their applications.

\SideHead{Folders}A folder icon.  Double-click it to open its window.

The name under each icon is the file's InitechDOS name.  Names longer than
the space under the icon are cut short on the screen, but not on the disk.

\section{Items that are not shown}
\index{hidden items}
Two items on the data disk are not shown in the \texttt{Drive A Files}
window, because the Finder uses them itself: the folder \file{TRASH},
which holds what you put in the Trash and is shown as the Trash window, and
the file \file{DESKTOP.DB}.  Do not change or delete them at the
InitechDOS prompt.

\section{Moving files and folders}
\index{moving files}
Dragging an icon to another folder moves it: the file or folder is taken
out of its folder and put into the other one, unchanged.  The Finder does
not copy files.  To copy a file, use the \file{COPY} or \file{XCOPY} command
at the InitechDOS prompt.

A drop is refused, and the icon slides back, when:
\begin{itemize}
\item the folder you drop it in already holds an item with the same name;
\item you drop a folder into itself or into one of its own folders;
\item you drop the icon back into the folder it came from;
\item you drop a locked item on the Trash (an alert says so).
\end{itemize}

\section{Selecting icons}
\index{selecting}
\begin{center}\small
\begin{tabular}{@{}p{46mm}p{53mm}@{}}
\textbf{To select} & \textbf{Do this}\\\hline
one icon & click it\\
several icons & click the first; hold down \key{Shift} and click the others\\
a group of icons & drag a rectangle round them from a bare part of the window\\
every icon in the window & choose \menuitem{Select All} from the \menuname{Edit} menu (\key{Ctrl-A})\\
nothing & click a bare part of the window or the desktop\\
\end{tabular}
\end{center}
